\documentclass[12pt]{article}

% --- Packages ---
\usepackage[utf8]{inputenc}
\usepackage[T1]{fontenc}
\usepackage[margin=1in]{geometry}
\usepackage{amsmath, amsfonts, amssymb, amsthm}
\usepackage{graphicx}
\usepackage{hyperref}
\usepackage{enumitem}
\usepackage{listings}
\usepackage{xcolor}
\usepackage{tcolorbox} % For the modern "Problem" boxes
\usepackage{fancyhdr}  % For headers and footers

% --- Visual Styling ---
\definecolor{courseblue}{RGB}{0, 51, 102}
\hypersetup{
    colorlinks=true,
    linkcolor=courseblue,
    urlcolor=blue,
}

% Code snippet styling
\definecolor{codegreen}{rgb}{0,0.6,0}
\definecolor{codegray}{rgb}{0.5,0.5,0.5}
\lstset{
    commentstyle=\color{codegreen},
    keywordstyle=\color{magenta},
    numberstyle=\tiny\color{codegray},
    stringstyle=\color{purple},
    basicstyle=\ttfamily\small,
    breaklines=true,
    frame=single,
    language=Python
}

% Custom Problem Box
\newtcolorbox{problem}[1]{
  colback=blue!5!white,
  colframe=courseblue,
  fonttitle=\bfseries,
  title=Problem #1
}

% --- Assignment Metadata ---
\newcommand{\courseName}{Temas selectos de Matemáticas II}
\newcommand{\courseURL}{https://eduenez.github.io/MathAIspring2026UTSA}
\newcommand{\assignmentTitle}{1st Assignment}
\newcommand{\studentName}{Navarrete Tepozán Iván} 

% --- Header/Footer Setup ---
\pagestyle{fancy}
\fancyhf{}
\lhead{\courseName}
\rhead{\studentName}
\cfoot{\thepage}


% --- Document Starts Here ---
\begin{document}

\begin{center}
    {\LARGE \textbf{\assignmentTitle}} \\
    \vspace{5pt}
    {\large \href{\courseURL}{Course Site @GitHub}} \\
    \vspace{10pt}
    \textbf{Student:} \studentName \hfill \textbf{Due Date:} {Wed 4 February}\\
    \rule{\linewidth}{0.4pt}
\end{center}

\section*{Warm-up Task}
\begin{problem}{}
Read the Written Work Guidelines and documents on mathematical writings referenced therein[cite: 188]. Write a short paragraph ($\approx 5$ lines) reflecting on one major error present in your math writing hitherto, and outlining some steps you will take to improve your writing for this course[cite: 189].
\end{problem}

\begin{proof}[Reflection]
% Escribe tu reflexión aquí.
There is one major mistake that I have been making over the years: sometimes I refuse to write a more elaborate explanation of the problem, because I assume that anyone reading my demonstration must have an equivalent level of mathematical knowledge. However, after reviewing the resource that was shared, I realize this is a mistake. A proof should be written so that it stands on its own, without relying on the reader to fill in the gaps.
\end{proof}

\section*{Task I}
Open the Jupyter (iPython) notebook \texttt{intro\_feedforward\_mnist.ipynb} (link opens the notebook in Google Colab) which creates and trains a simple feedforward-network to learn to recognize handwritten digits in the MNIST "database"[cite: 198]. Make sure to look at the MNIST Wikipedia page and Figure 1.9 in Deep Learning Ch. 1[cite: 199].

\begin{problem}{Sub-task 1}
Play around assigning different values to the parameters, (\texttt{dim\_hl\_1}, \texttt{dim\_hl\_2}) and the \texttt{num\_epochs}[cite: 205]. Make a table and provide a human-readable analysis of your findings[cite: 206]. 

For extra credit +20\%, also compare the 2-hidden layer feedforward NN to NNs with 1 or 3 hidden layers[cite: 207].
\end{problem}

\begin{proof}[Solution]
% Escribe tu solución y análisis aquí.
\end{proof}

\begin{problem}{Sub-task 2}
Choose one aspect or issue you would like to understand better, research it, and write about 1 page summarizing your findings[cite: 211]. This sub-task requires you to analyze and explain (not simply cut-and-paste content)[cite: 212]. Examples include (but are not limited to):
\begin{itemize}
    \item What is a TPU? What are its similarities and differences to a CPU and a GPU? [cite: 213]
    \item What is Keras? What is Jax? What is NumPy? How are they related? [cite: 214]
    \item What is a "logit"? What is a "ReLU"? What is "categorical cross-entropy"? [cite: 218]
    \item What are other datasets beyond MNIST broadly used as NN benchmarks? [cite: 219]
\end{itemize}
\end{problem}

\begin{proof}[Solution]
% Escribe tu investigación y resumen de 1 página aquí.
\end{proof}

\end{document}